% !TeX program = pdflatex
% =====================================================================
%  CS5055 - Internet of Things (2026-2)
%  Perfil de Proyecto Final - Descripcion breve
%  Compilar: pdflatex perfil_proyecto_CS5055.tex
% =====================================================================
\documentclass[10pt,a4paper]{article}

\usepackage[utf8]{inputenc}
\usepackage[T1]{fontenc}
\usepackage[spanish,es-noquoting]{babel}
\usepackage[left=1.5cm,right=1.5cm,top=1.2cm,bottom=1.5cm]{geometry}
\usepackage[scaled=0.92]{helvet}
\renewcommand{\familydefault}{\sfdefault}
\usepackage{booktabs}
\usepackage{array}
\usepackage{tabularx}
\usepackage{microtype}
\usepackage{xcolor}
\usepackage{enumitem}
\usepackage{titlesec}
\usepackage{fancyhdr}
\usepackage{hyperref}
\hypersetup{colorlinks=true,linkcolor=black,urlcolor=blue!55!black}

% --- Estilo de reporte: secciones compactas ---
\titleformat{\section}{\normalsize\bfseries}{\thesection.}{0.5em}{}
\titlespacing*{\section}{0pt}{8pt}{3pt}
\setlength{\parskip}{2.5pt}
\setlength{\parindent}{0pt}
\renewcommand{\arraystretch}{1.12}

\pagestyle{fancy}
\fancyhf{}
\renewcommand{\headrulewidth}{0pt}
\fancyfoot[C]{\scriptsize\color{black!55}CS5055 Internet of Things $\cdot$ Perfil de proyecto $\cdot$ \thepage}

\newcolumntype{L}[1]{>{\raggedright\arraybackslash}p{#1}}
\newcolumntype{C}[1]{>{\centering\arraybackslash}p{#1}}

% Rotulo de bloque, con aire propio arriba
\newcommand{\subhead}[1]{\vspace{6pt}\textbf{#1}\par\vspace{0pt}}

% Evitar lineas viudas y huerfanas
\clubpenalty=10000
\widowpenalty=10000

\begin{document}

% --- Encabezado tipo reporte ---
\noindent\rule{\textwidth}{1.2pt}\\[3pt]
{\large\bfseries Climate Cognitive System (CCS)}\\[2pt]
{\normalsize Nodo IoT de control térmico de aula: modelo de IA frente a heurística}\\[5pt]
{\footnotesize Curso: Internet of Things (CS5055) $\cdot$ UTEC $\cdot$ 2026-2}\\[1pt]
{\footnotesize Equipo: Brayan Eduardo Gomero Castillo $\cdot$ Kiara Alexandra Balcázar
Santa Cruz $\cdot$ Rodrigo Dion Carlos Acosta}\\[4pt]
\noindent\rule{\textwidth}{0.5pt}

% =====================================================================
\section{Problemática}

Las aulas de UTEC se climatizan con un sistema centralizado, donde la decisión de enfriar se
toma en un equipo central del edificio y no en cada aula. Esto genera tres problemas.

\begin{itemize}[leftmargin=1.2em,itemsep=1.5pt,topsep=2pt,parsep=0pt]
  \item \textbf{La climatización no responde a la ocupación.} El calor que se acumula en un
  aula depende de cuántas personas hay dentro, y esa cantidad cambia a lo largo del día.
  Como el sistema central no la mide, enfría aulas vacías y se queda corto cuando el aula
  está llena.
  \item \textbf{No se monitorea la calidad del aire.} El CO$_2$ aumenta conforme se llena el
  aula y afecta el confort y la concentración, pero hoy no se mide ni se corrige.
  \item \textbf{No existe alerta de monóxido de carbono.} Ante una fuga de CO no hay ningún
  sensor que la detecte ni ninguna respuesta automática.
\end{itemize}

El proyecto también responde una pregunta técnica: si un modelo de IA justifica su
complejidad frente a una heurística de umbrales para decidir cuándo climatizar. En lugar de
asumir la respuesta, la contrastamos con datos reales del aula.

% =====================================================================
\section{Sensores}

\vspace{2pt}
{\footnotesize
\begin{tabularx}{\textwidth}{@{}C{0.7cm} l X C{1.9cm}@{}}
\toprule
\textbf{\#} & \textbf{Sensor} & \textbf{Qué mide y para qué} & \textbf{Carácter} \\
\midrule
1 & DHT11 & Temperatura y humedad del aula, que son la entrada principal de las dos
    lógicas de decisión. & Vital \\
2 & MQ-135 & Calidad del aire, que usamos como estimación indirecta del CO$_2$. Sirve para
    calcular la ocupación y para activar la ventilación cuando el aire se vicia. & Vital \\
3 & MQ-7 & Monóxido de carbono, con el $R_0$ calibrado en aire limpio. Es el canal de
    seguridad del sistema. & Vital \\
4 & INA219 / ACS712 & Corriente y potencia que consume el actuador. Con esta medición el
    reporte de ahorro energético deja de ser una estimación. & Adicional \\
5 & Cámara + MediaPipe & Conteo de personas por cruce de línea, procesado en un equipo
    aparte del nodo. Valida la ocupación calculada a partir del CO$_2$. & Externo \\
\bottomrule
\end{tabularx}}

\vspace{3pt}
El mínimo de tres sensores que exige el curso lo cubren los tres primeros, que van
embarcados en el propio nodo ESP32. El sensor de consumo se declara como adicional porque
el sistema funciona igual sin él, y sin esa medición el ahorro energético se reporta como
estimación en lugar de dato medido. La cámara no se cuenta dentro de ese mínimo, porque su
procesamiento ocurre fuera del nodo, en un equipo aparte; queda como canal externo que
valida de forma independiente la ocupación estimada a partir del CO$_2$.

% =====================================================================
\section{Actuadores}

\vspace{2pt}
{\footnotesize
\begin{tabularx}{\textwidth}{@{}C{0.7cm} L{3.6cm} X C{1.9cm}@{}}
\toprule
\textbf{\#} & \textbf{Actuador} & \textbf{Función y disparo} & \textbf{Carácter} \\
\midrule
1 & Ventilador DC 12\,V con relé o MOSFET
  & Enfría y renueva el aire. Recibe la decisión por MQTT y la aplica encendiéndose y
    apagándose, o regulando el caudal con PWM. & Vital \\
2 & Buzzer + LED RGB
  & Alarma de CO. Se activa en el propio nodo al superar los 50\,ppm, sin depender de la
    red, y lleva el ventilador a caudal máximo. & Vital \\
3 & Calefactor PTC cerámico 12\,V
  & Calienta el aula. Añade la opción de calentar, de modo que la decisión puede ser
    calentar, enfriar o mantener. & Adicional \\
\bottomrule
\end{tabularx}}

\vspace{3pt}
Los dos primeros cumplen el requisito del curso y son independientes entre sí, ya que uno
responde a la temperatura y el otro a la seguridad. El calefactor es adicional. Todo el
prototipo funciona en baja tensión DC de hasta 12\,V y en ningún caso maneja
\mbox{220\,V AC}.

% =====================================================================
\section{Núcleo del proyecto: dos lógicas de decisión en paralelo}

Las dos lógicas trabajan sobre los mismos datos, y cada una decide por separado qué hacer
con el actuador.

\begin{itemize}[leftmargin=1.2em,itemsep=1.5pt,topsep=2pt,parsep=0pt]
  \item \textbf{Heurística.} Reglas de umbral sobre temperatura, CO$_2$ y ocupación. Incluye
  un margen de $\pm 0{,}5\,^\circ$C y un encendido mínimo de 60\,s para que el ventilador no
  se encienda y se apague de forma continua. Se puede revisar regla por regla y sirve como
  referencia de comparación.
  \item \textbf{Modelo de IA.} Una regresión con \emph{gradient boosting} que predice el
  calor que va a acumular el aula a partir de la temperatura, la hora, la ocupación esperada
  según la reserva, la ocupación real, la temperatura exterior y los datos del aula. Se
  valida con datos posteriores a los del entrenamiento y no se despliega si no supera a la
  heurística.
\end{itemize}

Las dos funcionan de forma permanente: una controla el ventilador y la otra decide en
paralelo sin actuar. El dashboard muestra ambas decisiones lado a lado con el motivo de cada
una, y el sistema registra cuánto coinciden, en qué casos difieren y cuánta energía consume
cada estrategia. Así el entregable no es solo un sistema que funciona, sino la evidencia de
cuál de las dos decide mejor. Si el modelo no supera a la heurística, ese resultado también
es válido y se reporta.

% =====================================================================
\section{Alcance}

\subhead{Qué se entrega.}
\begin{enumerate}[leftmargin=1.4em,itemsep=1.5pt,topsep=2pt,parsep=0pt]
  \item \textbf{Nodo IoT autónomo.} Un ESP32 con los sensores listados, PCB propia diseñada
  en KiCad y carcasa impresa en 3D en PLA, con la zona de sensado separada del flujo del
  ventilador para no alterar las lecturas.
  \item \textbf{Comunicación inalámbrica.} MQTT sobre WiFi en ambas direcciones con QoS 1,
  aviso automático cuando el nodo se desconecta, y el último comando queda guardado en el
  servidor para que el nodo recupere su estado tras un reinicio.
  \item \textbf{Lazo de control cerrado.} El recorrido completo, del sensado a la ocupación
  estimada, de ahí a la decisión de las dos lógicas, luego el comando por MQTT y finalmente
  la actuación física, con una latencia objetivo máxima de 3\,s.
  \item \textbf{Backend y dashboard web.} Incluyen procesamiento previo de los datos, con la
  predicción del calor que se va a acumular, la ocupación estimada, la comparación de las
  dos lógicas,
  las alertas de CO y el reporte de ahorro en kWh y USD sobre una ventana de 7 días.
  \item \textbf{Seguridad en el propio nodo.} La alarma de CO no depende del servidor, y si
  se pierde la comunicación el actuador pasa a un estado seguro.
  \item \textbf{Validación experimental.} El contraste entre la IA y la heurística con datos
  reales del aula, con al menos una figura de resultados propia.
\end{enumerate}

\subhead{Despliegue.}
Un nodo físico en un aula piloto, la A403 de UTEC Barranco. Añadir más aulas no requiere
cambios en el código, porque cada una envía y recibe sus mensajes por canales MQTT propios
identificados con su nombre, como \texttt{ccs/A403/telemetry}. Aun así, solo se demuestra un
nodo real y las aulas restantes se simulan.

\subhead{Fuera de alcance.}
\begin{itemize}[leftmargin=1.2em,itemsep=1.5pt,topsep=2pt,parsep=0pt]
  \item \textbf{El aire acondicionado central del campus.} No lo controlamos porque no ofrece
  ninguna forma de manejarlo desde el aula, ni por control remoto infrarrojo ni con un ajuste
  de temperatura local. Su operación está a cargo del personal del edificio y no es de acceso
  público, así que intervenirlo excede lo que corresponde a un proyecto de curso. Por eso el
  actuador del proyecto es local, y la integración con el sistema del campus queda como
  trabajo futuro.
  \item \textbf{Aplicativo móvil nativo.} No se implementa, y además es opcional según el
  enunciado. El dashboard web se adapta a pantallas de celular y se usa desde el navegador.
  \item \textbf{Tensión de red.} No se maneja 220\,V AC, ya que todo el prototipo trabaja en
  baja tensión DC.
  \item \textbf{Control de acceso.} El conteo de personas solo alimenta el modelo térmico, y
  ningún elemento del sistema abre, cierra ni bloquea un paso físico.
  \item \textbf{Precisión del MQ-135.} Lo usamos como una estimación relativa de CO$_2$ y no
  como una medición exacta en ppm. Se calibra en dos puntos, con el aula vacía y con el aula
  llena, y se contrasta con el conteo de la cámara.
\end{itemize}

\end{document}
